% Sección 2: resultados por fase, con evidencias y estado de cada uno

\seccion{Resultados por fase}

% ===============================================================
% 5. Fase 1: selección del conjunto de datos
% ===============================================================
\begin{frame}{Fase 1: selección del conjunto de datos}
    % [c]: las dos columnas centradas en vertical entre sí.
    % Centrado horizontal: columnas 0,44 + 0,38 con 0,04 entre ellas = 0,86.
    % Sobra 0,14, así que se corre la mitad (0,07) y totalwidth = 0,07 + 0,86 = 0,93.
    % Se deja 0,95 porque Beamer suma un pequeño espacio entre elementos
    \begin{columns}[c, totalwidth=0.95\textwidth]
        \hspace*{0.07\textwidth}
        \begin{column}{0.44\textwidth}
            Se seleccionó \textbf{CornOrb}, un conjunto de datos público:
            \begin{itemize}
                \item 1.454 ojos de 744 pacientes.
                \item 5.628 imágenes de mapas corneales.
                \item Cada ojo trae sus mapas y sus datos clínicos juntos, que es lo que necesita un modelo multimodal.
                \item El diagnóstico de cada ojo lo asignaron oftalmólogos.
            \end{itemize}
        \end{column}
        \hspace{0.04\textwidth}
        \begin{column}{0.38\textwidth}
            \centering
            {\small\textbf{Qué trae cada ojo}}\par\vspace{0.2cm}
            \begin{tikzpicture}[node distance=0.35cm]
                \node[caja=4.6cm] (m) {\textbf{4 mapas de la córnea}\\ \footnotesize Curvatura, elevación delantera, elevación trasera y espesor};
                \node[caja=4.6cm, below=of m] (v) {\textbf{Datos clínicos}\\ \footnotesize Edad, sexo, curvatura máxima, espesor y astigmatismo};
                \node[caja=4.6cm, draw=verde, below=of v] (e) {\textbf{Diagnóstico}\\ \footnotesize Normal o con queratocono};
            \end{tikzpicture}
        \end{column}
    \end{columns}
    \alcanzado{Objetivo 1}
\end{frame}

% ===============================================================
% 6. Fase 1: análisis exploratorio (composición y calidad)
% ===============================================================
\begin{frame}{Fase 1: revisión de los datos}
    \begin{columns}[T]
        \begin{column}{0.52\textwidth}
            % Exportada por el EDA al ancho de esta columna: se inserta sin escalar
            \centering
            \includegraphics[width=\linewidth]{eda/eda_clases.pdf}
        \end{column}
        \begin{column}{0.44\textwidth}
            \vspace{0.5 cm}
            \footnotesize
            \textbf{Puntos positivos}
            \begin{itemize}
                \item No faltan datos ni imágenes de ningún paciente.
                \item Los ojos de un mismo paciente nunca quedan repartidos entre entrenamiento y prueba.
            \end{itemize}
            \textbf{Puntos a tratar}
            \begin{itemize}
                \item 46 ojos medidos más de una vez.
                \item Un código de paciente mal escrito.
                \item {Valores imposibles, como una curvatura\\de 4 cuando lo normal es 44.}
            \end{itemize}
        \end{column}
    \end{columns}
    \alcanzado{Objetivo 1}
\end{frame}

% ===============================================================
% 7. Fase 1: qué dicen las variables clínicas
% ===============================================================
\begin{frame}{Fase 1: qué dicen los datos clínicos}
    \begin{columns}[T]
        \begin{column}{0.56\textwidth}
            \vspace{-0.25 cm}
            % Exportada por el EDA al ancho de esta columna: se inserta sin escalar
            \centering
            \includegraphics[width=\linewidth]{eda/eda_correlacion.pdf}
        \end{column}
        \begin{column}{0.40\textwidth}
            \vspace{0.5 cm}
            \footnotesize
            Los ojos con queratocono tienen:
            \begin{itemize}
                \item {Una córnea más curva\\(53 frente a 44 dioptrías)}.
                \item {Una córnea más delgada\\(428 frente a 521 micrones)}.
                \item Más astigmatismo.
            \end{itemize}
            Coincide con lo que describe la medicina sobre la enfermedad.
        \end{column}
    \end{columns}
    \idea{Los datos clínicos ya separan bastante los dos grupos.\\Por eso el modelo los combina con las imágenes.}
\end{frame}

% ===============================================================
% 8. Datos locales para la validación
% ===============================================================
\begin{frame}{Datos locales para la validación}
    % Texto más ancho y bloques cerca: 0,56 + 0,03 + 0,34 = 0,93.
    % Sobra 0,07: se corre la mitad (0,035) y totalwidth queda con holgura en 0,98
    \begin{columns}[c, totalwidth=0.98\textwidth]
        \hspace*{0.035\textwidth}
        \begin{column}{0.56\textwidth}
            CornOrb viene de pacientes de Argelia. Para comprobar que KERALYTIX funciona con pacientes de la región, hacen falta estudios locales.
            \vspace{-0.5 cm}
            \begin{itemize}
                \item Se contactaron 9 centros y profesionales de Misiones, Córdoba, Buenos Aires y Perú.
                \item Se contactó, por medio del programador Gonzalo Pozzo, con el expresidente de la Sociedad Internacional de Queratocono Roberto Albertazzi. 
            \end{itemize}
        \end{column}
        \hspace{0.03\textwidth}
        \begin{column}{0.34\textwidth}
            \centering
            \begin{tikzpicture}[node distance=0.4cm]
                \node[caja=4.2cm, draw=verde] (a) {Contacto con Roberto Albertazzi};
                \node[caja=4.2cm, fill=azul!15, very thick, below=of a] (c) {Trámites legales para la solicitud\\ \footnotesize estamos acá};
                \node[caja=4.2cm, draw=gray!60, text=gray, below=of c] (d) {Estudios locales anonimizados};
                \draw[flecha] (a) -- (c);
                \draw[flecha, gray!60] (c) -- (d);
            \end{tikzpicture}
        \end{column}
    \end{columns}
    \enproceso{Validación local}
\end{frame}

% ===============================================================
% 9. Fase 2: cómo funciona el modelo (diagrama de bloques)
% ===============================================================
\begin{frame}{Fase 2: cómo funciona el modelo}
    \begin{center}
        \begin{tikzpicture}[node distance=0.5cm and 0.55cm]
            \node[caja=1.9cm] (mapas) {\textbf{Mapas}};
            \node[caja=1.9cm, below=1.1cm of mapas] (datos) {\textbf{Datos clínicos}};
            \node[caja=3cm, right=of mapas] (cnn) {\textbf{Rama de imágenes}\\ \footnotesize aprende a reconocer formas en los mapas};
            \node[caja=3cm, right=of datos] (densa) {\textbf{Rama de datos}\\ \footnotesize analiza los valores numéricos};
            % Todas las flechas miden 0,55 cm (la separación horizontal de node distance).
            % Las dos ramas se unen en (j), a 0,55 cm, y de ahí sale una sola flecha
            % del doble de largo (1,1 cm) hasta la fusión
            \coordinate (m2) at ($(cnn.east)!0.5!(densa.east)$);
            \coordinate (j) at ($(m2)+(0.55cm,0)$);
            \node[caja=2cm, fill=azul!15, very thick, right=0.45cm of j] (fus) {\textbf{Fusión}\\ \footnotesize junta lo que encontró cada rama};
            \node[caja=2.6cm, draw=verde, right=of fus] (res) {\textbf{Resultado}\\ \footnotesize normal o queratocono, y confianza};
            \node[caja=2.6cm, draw=verde, above=0.45cm of res] (gc) {\textbf{Mapa de calor}\\ \footnotesize zonas que influyeron};
            \draw[flecha] (mapas) -- (cnn);
            \draw[flecha] (datos) -- (densa);
            \draw[flecha] (cnn.east) -- ++(0.35cm,0) |- (fus.west);
            \draw[flecha] (densa.east) -- ++(0.35cm,0) |- (fus.west);
            \draw[flecha] (fus) -- (res);
            \draw[flecha] (cnn.north) |- node[pos=0.75, above, font=\scriptsize, text=gray] {Grad-CAM} (gc.west);
        \end{tikzpicture}
    \end{center}
    {\small Cada rama se especializa en un tipo de dato y la fusión aprovecha las relaciones entre ambos. Grad-CAM mira la rama de imágenes y marca qué zonas pesaron en el resultado.}
\end{frame}

% ===============================================================
% 10. Fase 2: primeros resultados del modelo de imágenes
% ===============================================================
\begin{frame}{Fase 2: primeros resultados}
    % Columnas juntas y centradas: 0,5 + 0,03 + 0,34 = 0,87.
    % Sobra 0,13: se corre la mitad (0,065) y totalwidth queda con holgura en 0,96
    \begin{columns}[c, totalwidth=0.96\textwidth]
        % \hspace*{0.065\textwidth}
        \begin{column}{0.6\textwidth}
            \small
            Se compararon tres \textbf{redes base}, modelos ya entrenados con millones de imágenes que se adaptan a los mapas corneales. En esta primera prueba la que mejor funcionó fue MobileNetV2, la más liviana. Todavía se siguen probando otras redes y estrategias.
            \vspace{0.2cm}

            Resultado de prueba con 279 ojos que el modelo nunca vio:
            \begin{itemize}
                \item Detectó \textbf{95 de 103} ojos con queratocono (92\,\%).
                \item \textbf{Ninguna falsa alarma} en los 176 ojos sanos.
            \end{itemize}
            \vspace{0.1cm}
            Un modelo solo con los datos clínicos da resultados parecidos. Por eso el siguiente paso es combinar ambos.
        \end{column}
        \hspace{0.03\textwidth}
        \begin{column}{0.34\textwidth}
            \centering
            \begin{tikzpicture}[node distance=0.3cm]
                \node[caja=4.2cm, draw=verde] (a) {Datos limpios y separados por paciente};
                \node[caja=4.2cm, draw=verde, below=of a] (b) {Primera prueba con los mapas corneales};
                \node[caja=4.2cm, fill=azul!15, very thick, below=of b] (c) {\textbf{Pruebas de otras redes y estrategias}\\ \footnotesize estamos acá};
                \node[caja=4.2cm, draw=gray!60, text=gray, below=of c] (d) {Fusión con los datos clínicos};
                \draw[flecha] (a) -- (b);
                \draw[flecha] (b) -- (c);
                \draw[flecha, gray!60] (c) -- (d);
            \end{tikzpicture}
        \end{column}
    \end{columns}
    \enproceso{Objetivo 2}
\end{frame}

% ===============================================================
% 11. Fase 2: explicación visual con Grad-CAM
% ===============================================================
\begin{frame}{Fase 2: el modelo muestra dónde miró}
    % Columna de imágenes justa para la grilla (0,45) y texto pegado a ella:
    % 0,45 + 0,03 + 0,42 = 0,90. Se corre la mitad de lo que sobra (0,05)
    % y totalwidth queda con holgura en 0,97
    \begin{columns}[c, totalwidth=0.97\textwidth]
        \hspace*{0.05\textwidth}
        \begin{column}{0.45\textwidth}
            \centering
            \begin{tabular}{@{}c@{\hspace{0.15cm}}c@{\hspace{0.15cm}}c@{}}
                & {\scriptsize Mapa de la córnea} & {\scriptsize Mapa de calor} \\
                \rotatebox{90}{\scriptsize\hspace{0.4cm}Con queratocono} &
                \includegraphics[width=0.44\linewidth]{gradcam/qc_mapa.png} &
                \includegraphics[width=0.44\linewidth]{gradcam/qc_gc.png} \\
                \rotatebox{90}{\scriptsize\hspace{0.9cm}Sano} &
                \includegraphics[width=0.44\linewidth]{gradcam/sano_mapa.png} &
                \includegraphics[width=0.44\linewidth]{gradcam/sano_gc.png} \\
            \end{tabular}
        \end{column}
        \hspace{0.03\textwidth}
        \begin{column}{0.42\textwidth}
            \small
            Mapas de la cara posterior de dos ojos de prueba.
            \begin{itemize}
                \item En el ojo con queratocono, el calor cae sobre la zona deformada. El modelo lo dio como queratocono con total confianza.
                \item En el ojo sano no marca la córnea central y lo dio como sano.
            \end{itemize}
            Así el oftalmólogo puede comprobar en qué se basó cada resultado.
        \end{column}
    \end{columns}
    \enproceso{Objetivo 3}
\end{frame}

% ===============================================================
% 11. Fase 3: la aplicación web
% ===============================================================
\begin{frame}{Fase 3: la aplicación web}
    \begin{center}
        \begin{tikzpicture}[node distance=0.5cm and 0.65cm]
            \node[caja=2.3cm] (med) {\textbf{Oftalmólogo}\\ \footnotesize desde el navegador};
            \node[caja=2.8cm, right=of med] (front) {\textbf{Página web}\\ \footnotesize carga los estudios y muestra el informe};
            \node[caja=2.8cm, right=of front] (back) {\textbf{Servidor}\\ \footnotesize recibe los estudios y ejecuta el modelo};
            \node[caja=1.9cm, fill=azul!15, very thick, right=of back] (mod) {\textbf{Modelo de IA}};
            % Nota debajo de la página web, sin recuadro
            \node[below=0.15cm of front, font=\footnotesize, text=gray] {Modelo de suscripción};
            \draw[flecha, <->] (med) -- (front);
            \draw[flecha, <->] (front) -- (back);
            \draw[flecha, <->] (back) -- (mod);
        \end{tikzpicture}
    \end{center}
    \vspace{0.2cm}
    % Abajo, dos columnas: el texto a la izquierda y el QR al HTML de Fabricio a la derecha
    \begin{columns}[c, totalwidth=0.97\textwidth]
        \hspace*{0.05\textwidth}
        \begin{column}{0.6\textwidth}
            \begin{itemize}
                \small
                \item La página web se hace con React y el servidor con FastAPI.
                \item Se trabaja en noviembre y diciembre, superpuesta con el cierre del modelo.
            \end{itemize}
        \end{column}
        \hspace{0.03\textwidth}
        \begin{column}{0.25\textwidth}
            \centering
            % Reemplazar el recuadro por el QR, por ejemplo:
            % \includegraphics[height=2.2cm]{qr_web.png}\par {\scriptsize Escaneá para ver la página}
            \fcolorbox{gray!60}{white}{\parbox[c][2.2cm][c]{2.2cm}{\centering\scriptsize\color{gray} QR de la\\ página web}}
        \end{column}
    \end{columns}
    \porhacer{Objetivo 4}
\end{frame}
